\begin{afterword}
\summaryhead

The essay argues that human feelings of care do not scale with numbers. A million and a billion feel much the same to the author, whose ``care-o-meter'' cannot register the suffering of billions of people or the stakes of the far future. So feelings cannot guide action on large problems.

Three invented donors, Alice, Bob and Christine, give small sums because of social pressure, guilt or competition. A fourth, Daniel, imagines saving one oiled bird. He decides it is worth three minutes or \$3 to him, and multiplies by the number of oiled birds. He concludes that he ``actually'' cares more than he can afford, and then that every problem is his, with too many problems to address.

The author says this is not preaching. Prominent altruists, in the author's account, are people who have learned not to trust their feelings. The essay recommends doing the multiplication, trusting numbers over feelings, and acting with desperation but without guilt, and promotes Giving What We Can, GiveWell, MIRI and the Future of Humanity Institute.

\responsehead

The essay makes one observation and draws one rule from it. The observation is sound: feelings of concern do not grow in proportion to the number of people affected. The rule is to put a number on what one person or one bird is worth, multiply, and trust the product over the feeling. The essay tells this as a personal shift, and it ends by inviting the reader to make the same shift and ``join the ranks of people saving the world.''

The rule depends on an assumption the essay never states: that each bird is worth the same however many birds there are. That is the whole content of ``multiply,'' and it is contestable. One can coherently value the thousandth bird less than the first, or care about a population and not about each of its members. The essay does not argue against these views. It treats the scaled-up feeling as the true one.

The one worked example does not end in any action. Daniel does the sum, feels ``growing horror,'' and gives neither to the birds nor to the two other charities named. What changes is that he no longer finds the question rude. The only action is a hope, in a parenthesis, that he will go on to discover effective altruism. For poverty and the far future, where the essay says the numbers call for more resources than exist, no multiplication is shown.

The essay also uses an example that works against it. Gandhi, Mother Teresa and Nelson Mandela are offered as altruists who did not simply care more but learned not to trust their feelings. Paul Slovic, whose research documents the effect this essay describes, took the title of a 2007 paper from Mother Teresa: ``If I look at the mass I will never act. If I look at the one, I will.'' That is the opposite habit: acting on the feeling for one person. Slovic's own conclusion from the same diagnosis was to rely on argument and law, not on private desperation.

The list of organizations ``to start'' includes MIRI, where the author had joined as a research fellow earlier that year (a 2015 announcement by the author puts the time there at ``a little over a year''). The essay does not say so. The list runs from poverty charities to the long-term future, and the only link between them in the text is a single unsourced figure of ``quadrillions'' of possible future beings.

In short: a sound observation about feeling, turned into a rule that assumes value adds up bird by bird, shown in one example that ends without action, and closed with an undisclosed recommendation of the author's employer.
\end{afterword}
